\scanpage{19}
\mspage{18}

la chose essentielle est la connaissance de
\[
\underset{\unsure{M|m}}{\mathrel{\tikz[baseline=-.55ex,x=1em,y=1em]{\draw[line width=.065em,line cap=round,line join=round] (0,.40)--(.48,0)--(0,-.35) (.51,.40)--(.03,.02)--(.49,-.35);}}}\quad\text{relation binaire sym. réfl. (sur }M\text{)}
\tag*{(1.35)}
\]
\deleted{\unsure{qui indique}} satisfaisant
\[
 \unsure{X|x} \mathrel{\tikz[baseline=-.55ex,x=1em,y=1em]{\draw[line width=.065em,line cap=round,line join=round] (0,.40)--(.48,0)--(0,-.35) (.51,.40)--(.03,.02)--(.49,-.35);}} \unsure{Y|y},\quad X'\leq X,\,Y'\leq Y\quad\Longrightarrow\quad X' \mathrel{\tikz[baseline=-.55ex,x=1em,y=1em]{\draw[line width=.065em,line cap=round,line join=round] (0,.40)--(.48,0)--(0,-.35) (.51,.40)--(.03,.02)--(.49,-.35);}} \unsure{Y'|Y,}
 \tag*{(1.35)}
\]

\deleted{\unsure{elle|alors}}

il y a corr. biuniv. \unsure{entre|en} l'ens \unsure{des|de} telles relations sur $M$, et l'ens des $\unsure{\mathcal{F}|\mathscr{F}}\inserted{\subset\unsure{\mathcal{P}|P}(M)}$ satisfaisant les conditions énumérées précédemment, plus celle\deleted{\unsure{s} d'être} que les $F\in\mathcal{F}$ soient \deleted{\unsure{toutes|toute}} de type fini\unsure{.|e} C'est en ce sens que j'ai pu dire que ``\unsure{moralement|essentiellement}'', le couple $(M,\mathcal{F}\subset\mathcal{P}(M))$ équivalait à \deleted{\illegible{1}} la donnée de $\underset{M}{\leqslant}$ \unsure{seul|et|soit} \unsure{de|des} $\mathrel{\tikz[baseline=-.55ex,x=1em,y=1em]{\draw[line width=.065em,line cap=round,line join=round] (0,.40)--(.48,0)--(0,-.35) (.51,.40)--(.03,.02)--(.49,-.35);}}$ satisfaisant (1.36).

\deleted{Mais il s'agit maintenant}

Pratiquement, dans presque tous les exemples, c'est via $M$\inserted{(\unsure{multistrates|multistrate})}, muni de $\leq$ et d'une relation de disjonction $\mathrel{\tikz[baseline=-.55ex,x=1em,y=1em]{\draw[line width=.065em,line cap=round,line join=round] (0,.40)--(.48,0)--(0,-.35) (.51,.40)--(.03,.02)--(.49,-.35);}}$, qu'on \unsure{définira|définirait} les ateliers\ldots

Il faut à présent en venir à la relation $\ll$. \deleted{\unsure{Comme}} \unsure{Elle} va pour l'instant \unsure{servir} surtout à former des ombres\unsure{,|.} Je commence par les grandes ombres
